\subsection{Scope and the precise unresolved question}
We fix a universe and take all categories to be small in it. We use
covariant functors $H:A\to\Set$, with arrows in $\El_AH$ directed by
$(a,h)\xrightarrow{f}(b,H(f)h)$. We put $BA=|NA|$. A nonempty connected
category $A$ is total if each component of every such category of
elements is weakly contractible. Exactness means exactness of
$\colim_A:\Ab^A\to\Ab$. The module $P_p=\Z[A(p,-)]$ is covariant;
$\Delta\Z$ is the constant right $\Z A$-module. We write
$T_i(M)=\Tor_i^{\Z A}(\Delta\Z,M)=L_i\colim_A M$.

\begin{obstruction}[Open in this incomplete record: Part A.1]\label{A:main}
We have not established whether every total category is exact. Equivalently,
we have not proved that every finite parallel family
$a_0,\ldots,a_k:p\to q$ has an integral affine post-equalizer.
\end{obstruction}
For $k=1$ the assertion follows from the audited permutation-module
calculation in Lemma~\ref{B:row-audit}. For $k\geq2$, the transposed
presentation of Proposition~\ref{B:lead3} detects the required
simultaneous annihilator. None of the calculations of the row
presentation identifies its kernel with this intersection. An affirmative
answer for triples alone would still leave larger finite families:
postcomposing by an affine equalizer creates affine arrows, to which the
ordinary triple hypothesis does not directly apply.

\subsection{An arithmetic reduction}
\begin{proposition}[Proved: fieldwise detection of integral exactness]\label{A:fieldwise}
For every small category $A$, including the empty category and categories
with several components, the following conditions are equivalent:
\begin{enumerate}
\item $\colim_A:\Ab^A\to\Ab$ is exact;
\item $\colim_A:(\Q\text{-}\mathrm{Mod})^A\to\Q\text{-}\mathrm{Mod}$
is exact and, for every prime $p$, the functor
$\colim_A:(\F_p\text{-}\mathrm{Mod})^A\to\F_p\text{-}\mathrm{Mod}$
is exact.
\end{enumerate}
No totality hypothesis is needed. Consequently, Parts A.1 and A.2 taken
together ask equivalent universal questions. A negative answer to A.1
necessarily gives a failure over $\Q$ or over at least one prime field,
although the obstructing diagram need not remain the same after reduction.
\end{proposition}
\begin{proof}
The ordinary abelian-group colimit of a diagram of $R$-modules, for
$R=\Q$ or $\F_p$, carries the unique compatible $R$-module structure,
and its underlying abelian group is the ordinary abelian-group colimit.
This follows from the direct-sum-and-relations construction, since the
subgroup of relations is an $R$-submodule. Forgetting scalars is exact,
so (1) implies (2).

Assume (2). We use the normalized bar complex
\[
 C_n(A;M)=\bigoplus_{a_0\to\cdots\to a_n}^{\mathrm{nondegenerate}}
 M(a_0),
\]
with its usual face differential; its homology is $T_n(M)$. If $M$ is
a $\Q$- or $\F_p$-module diagram, this is, after restriction of scalars,
exactly the bar complex computing derived colimits over that field.
Thus such diagrams have $T_i(M)=0$ for $i>0$. This comparison is an
equality of the specified complexes, not an assumption that arbitrary
base change is flat.

We recall the three closure facts used here. A direct sum or filtered
colimit of acyclic diagrams is acyclic because the bar complex commutes
with these operations and both operations are exact in $\Ab$. If
$0\to U\to V\to W\to0$ is objectwise exact, the bar complexes form a
short exact sequence; their long exact homology sequence shows closure
under extensions and under kernels of epimorphisms between acyclic
diagrams.

Let $M:A\to\Ab$ be arbitrary. For each prime $p$, its $p$-primary
torsion subfunctor is the union of $M[p^n]$, where
$M[p^n](a)=\{m:p^nm=0\}$. The quotient
$M[p^n]/M[p^{n-1}]$ is an $\F_p$-module diagram. Induction on $n$,
then filtered colimits, prove acyclicity of the $p$-primary subfunctor.
The full torsion subfunctor $M_{\mathrm{tor}}$ is their direct sum and
is therefore acyclic. Put $M_{\mathrm{tf}}=M/M_{\mathrm{tor}}$.
It is objectwise torsion free, so the natural map
\[
 M_{\mathrm{tf}}\longrightarrow\Q\otimes_\Z M_{\mathrm{tf}}
\]
is a monomorphism in $\Ab^A$. Its cokernel is objectwise torsion:
each element is a finite rational linear combination, and a common
denominator kills its class. Both the middle diagram and this cokernel
are acyclic. Closure under kernels of epimorphisms makes
$M_{\mathrm{tf}}$ acyclic. Finally the sequence
$0\to M_{\mathrm{tor}}\to M\to M_{\mathrm{tf}}\to0$ makes $M$
acyclic. Therefore $T_1(M)=0$ for every $M$, which is equivalent to
exactness of the right-exact colimit functor.
\end{proof}

\begin{obstruction}[Open in this incomplete record: Part A.2]\label{A:coefficients}
We have not established rational exactness for every total category, nor
exactness over each $\F_p$. Proposition~\ref{A:fieldwise} is a reduction,
not a decision of either implication.
\end{obstruction}

\subsection{Proof and construction routes investigated}
The accompanying report records the precise scope of each attempted
route. The model-category observation that totality makes ordinary
simplicial-set colimits homotopy invariant does not by itself identify
the underlying set colimit with the abelian-group colimit. A free
abelian cotriple resolution has permutation modules as its terms, but
some face maps send a basis vector to a linear combination; it is not
the linearization of a simplicial set diagram. We therefore do not infer
acyclicity of its colimit complex from totality. The row/column formulas
below isolate the unresolved first simultaneous-intersection instance.
